\section{Limitations and Conclusion}
\label{sec:conclusion}

\paragraph{Limitations.} The benchmark is simulator-based. Its ground truth is internally consistent, but physical fidelity and visual realism are bounded by the engine and renderer. The multiple-choice format gives a clean behavioral signal, yet it does not test open-ended hypothesis writing or expose the mechanism a model used to answer. Its difficulty also depends on how close the distractors are to the correct future; the simulator-oracle margin bounds this dependence but does not remove it. Harness choices also matter; reporting model--harness systems makes this dependence explicit but does not exhaust the harness space. Finally, the latent-law space is sampled. Seven phenomenon families cover several common modes of physical variation, not all possible physical worlds. \TODO{Add result-specific limitations after the baseline analysis.}

\paragraph{Conclusion.} We introduced a benchmark that asks whether interaction helps a multimodal agent understand an unfamiliar physical world, and whether that understanding carries over from interaction to vision. Interaction and evaluation are the same visual experiment seen from two sides: the agent first poses experiments and observes their outcomes, then, with interaction frozen, foresees the outcomes of new experiments posed by the world. Online learning curves, the budget sweep, and cross-realization and canonical-physics controls separate local use of an interaction history from an understanding of the particular world's laws. \TODO{Insert the main empirical takeaway once results land.} Future work should broaden the phenomenon families and intervention vocabulary, pair the behavioral test with analyses that localize failures inside the model--harness system, and move from rendered scenes toward agents that learn physical worlds through interaction and apply that understanding to what they see in the real world.

\paragraph{Reproducibility.} \TODO{State the release plan for the simulator, renderer, world and scene generators, harness configurations, and evaluation code.}
